%! TEX TS-program = xelatex
% MarquezSalazarBrandon-Interpolaciones.tex - MAIN FILE

\documentclass[12pt, a4paper]{article}

% ==================== CONFIGURACIÓN BÁSICA ====================
\usepackage{geometry}
\geometry{
    margin=2.5cm,
    headheight=14.5pt,
    includeheadfoot
}

\usepackage{polyglossia}
\setmainlanguage{spanish}

% ==================== FUENTES ====================
\usepackage{fontspec}
\setmainfont{QTBookmann}
\setmonofont{VictorMono Nerd Font}

% ==================== PAQUETES ESENCIALES ====================
\usepackage{amsmath, amssymb, amsthm}
\usepackage{graphicx}
\usepackage{xcolor}
\usepackage{tikz}
\usetikzlibrary{shapes, arrows, positioning, matrix, calc}

% ==================== FORMATO DEL DOCUMENTO ====================
\usepackage{fancyhdr}
\pagestyle{fancy}
\fancyhf{}
\rhead{\textsc{Brandon Márquez Salazar}}
\lhead{\textsc{Interpolación Polinomial Bayer}}
\cfoot{\thepage}
\renewcommand{\headrulewidth}{0.4pt}

\usepackage{titlesec}
\titleformat{\section}{\normalfont\Large\bfseries\scshape}{\thesection}{1em}{}
\titleformat{\subsection}{\normalfont\large\bfseries}{\thesubsection}{1em}{}
\titleformat{\subsubsection}{\normalfont\normalsize\bfseries\itshape}{\thesubsubsection}{1em}{}

% ==================== LISTADOS Y CÓDIGO ====================
\usepackage{listings}
\usepackage{lstautogobble}

% Configuración para incrustar archivos de código
\lstset{
    basicstyle=\ttfamily\small,
    numbers=left,
    numberstyle=\tiny\color{gray},
    frame=single,
    breaklines=true,
    captionpos=b,
    autogobble=true,
    showstringspaces=false,
    tabsize=2
}

% Colores para los listados
\definecolor{codegreen}{rgb}{0,0.6,0}
\definecolor{codegray}{rgb}{0.5,0.5,0.5}
\definecolor{codepurple}{rgb}{0.58,0,0.82}
\definecolor{backcolour}{rgb}{0.95,0.95,0.92}

\lstdefinestyle{mystyle}{
    backgroundcolor=\color{backcolour},
    commentstyle=\color{codegreen},
    keywordstyle=\color{magenta},
    numberstyle=\tiny\color{codegray},
    stringstyle=\color{codepurple},
    basicstyle=\ttfamily\footnotesize,
    breakatwhitespace=false,
    breaklines=true,
    captionpos=b,
    keepspaces=true,
    numbers=left,
    numbersep=5pt,
    showspaces=false,
    showstringspaces=false,
    showtabs=false,
    tabsize=2
}
\lstset{style=mystyle}

% ==================== HIPERVÍNCULOS (ÚLTIMO) ====================
\usepackage{hyperref}
\hypersetup{
    colorlinks=true,
    linkcolor=blue,
    filecolor=magenta,
    urlcolor=cyan,
    pdftitle={Interpolación Polinomial Bayer - Brandon Márquez Salazar},
    pdfauthor={Brandon Márquez Salazar}
}
\usepackage{float}

\begin{document}

% ==================== PORTADA ====================
% portada.tex
\begin{titlepage}
    \centering
    \vspace*{2cm}
    
    {\huge\bfseries Interpolación Bicuadrada y Bicúbica \\ para Patrón Bayer RGGB}
    
    \vspace{1.5cm}
    
    {\Large Brandon Marquez Salazar}
    
    \vspace{0.5cm}
    
    {\large Procesamiento Digital de Señales en Tiempo Real\\
    Universidad de Guanajuato}
    
    \vspace{2cm}
    
    % Gráfico mejorado del patrón Bayer
    \begin{tikzpicture}[scale=1.2]
        % Cuadrícula principal
        \draw[black, thick] (0,0) rectangle (2,2);
        \draw[black, thick] (1,0) -- (1,2);
        \draw[black, thick] (0,1) -- (2,1);
        
        % Canales de color
        \fill[red!80!black] (0,1) rectangle (1,2);
        \node[white, font=\bfseries] at (0.5,1.5) {R};
        
        \fill[green!60!black] (1,1) rectangle (2,2);
        \node[white, font=\bfseries] at (1.5,1.5) {G₁};
        
        \fill[green!40!black] (0,0) rectangle (1,1);
        \node[white, font=\bfseries] at (0.5,0.5) {G₂};
        
        \fill[blue!80!black] (1,0) rectangle (2,1);
        \node[white, font=\bfseries] at (1.5,0.5) {B};
        
        % Etiquetas
        \node[below, font=\small] at (1,-0.2) {Patrón Bayer RGGB 2×2};
    \end{tikzpicture}
    
    \vfill
    
    {\large \today}
\end{titlepage}

\tableofcontents
\newpage

% ==================== CONTENIDO PRINCIPAL ====================
% introduccion.tex
\section{Introducción}

\subsection{Contexto del Problema}
Los sensores de imagen en cámaras digitales utilizan filtros de color para capturar información cromática. 
El patrón Bayer, desarrollado por Bryce Bayer en 1976, es el arreglo más común, donde cada fotosito del sensor 
captura únicamente un canal de color (rojo, verde o azul).

\subsection{Problema de Interpolación}
Dada una matriz 2×2 del patrón Bayer:

\[
\begin{bmatrix}
R & G \\
G & B
\end{bmatrix}
\]

Es necesario estimar los valores faltantes en cada posición para reconstruir tres canales completos (RGB) 
con la misma resolución espacial. Este proceso se denomina \textbf{demosaicing} o interpolación de Bayer.

\subsection{Objetivos del Estudio}
\begin{enumerate}
    \item Implementar algoritmos de interpolación bicuadrada para el patrón Bayer RGGB
    \item Implementar algoritmos de interpolación bicúbica para el patrón Bayer RGGB
    \item Comparar los resultados con métodos de interpolación lineal tradicional
    \item Analizar la complejidad computacional de cada método
    \item Visualizar los procesos mediante diagramas y gráficos explicativos
    \item Proporcionar implementaciones prácticas en MATLAB para uso educativo
\end{enumerate}

\subsection{Organización del Documento}
Este documento se estructura de la siguiente manera: la Sección \ref{sec:metodos} describe los métodos 
matemáticos utilizados; la Sección \ref{sec:resultados} presenta los resultados obtenidos; la Sección 
\ref{sec:discusion} analiza los hallazgos; y finalmente, la Sección \ref{sec:conclusion} ofrece conclusiones 
y trabajo futuro. Los archivos de código fuente se encuentran en el Apéndice.
% metodos.tex
\section{Métodos}\label{sec:metodos}

\subsection{Patrón Bayer RGGB}

El patrón Bayer RGGB organiza los filtros de color en una matriz donde:
\begin{itemize}
    \item Las filas pares comienzan con rojo (R) y verde (G) alternados
    \item Las filas impares comienzan con verde (G) y azul (B) alternados
\end{itemize}

\begin{figure}[h]
\centering
\begin{tikzpicture}[scale=0.8]
    % Cuadrícula 4x4
    \draw[gray!30, very thin] (0,0) grid (4,4);
    \draw[black, thick] (0,0) rectangle (4,4);
    
    % Fila 1: R G R G
    \fill[red] (0,3) rectangle (1,4);
    \fill[green!70!black] (1,3) rectangle (2,4);
    \fill[red] (2,3) rectangle (3,4);
    \fill[green!70!black] (3,3) rectangle (4,4);
    
    % Fila 2: G B G B
    \fill[green!50!black] (0,2) rectangle (1,3);
    \fill[blue] (1,2) rectangle (2,3);
    \fill[green!50!black] (2,2) rectangle (3,3);
    \fill[blue] (3,2) rectangle (4,3);
    
    % Fila 3: R G R G
    \fill[red] (0,1) rectangle (1,2);
    \fill[green!70!black] (1,1) rectangle (2,2);
    \fill[red] (2,1) rectangle (3,2);
    \fill[green!70!black] (3,1) rectangle (4,2);
    
    % Fila 4: G B G B
    \fill[green!50!black] (0,0) rectangle (1,1);
    \fill[blue] (1,0) rectangle (2,1);
    \fill[green!50!black] (2,0) rectangle (3,1);
    \fill[blue] (3,0) rectangle (4,1);
    
    % Etiquetas
    \node[white] at (0.5,3.5) {R};
    \node[white] at (1.5,3.5) {G};
    \node[white] at (2.5,3.5) {R};
    \node[white] at (3.5,3.5) {G};
    
    \node[white] at (0.5,2.5) {G};
    \node[white] at (1.5,2.5) {B};
    \node[white] at (2.5,2.5) {G};
    \node[white] at (3.5,2.5) {B};
    
    \node[white] at (0.5,1.5) {R};
    \node[white] at (1.5,1.5) {G};
    \node[white] at (2.5,1.5) {R};
    \node[white] at (3.5,1.5) {G};
    
    \node[white] at (0.5,0.5) {G};
    \node[white] at (1.5,0.5) {B};
    \node[white] at (2.5,0.5) {G};
    \node[white] at (3.5,0.5) {B};
    
    \node[below] at (2,-0.3) {Patrón Bayer RGGB 4×4};
\end{tikzpicture}
\caption{Distribución de canales en patrón Bayer RGGB}
\end{figure}

\subsection{Interpolación Lineal}
La interpolación lineal es el método más simple, usando el promedio de vecinos inmediatos:

\[
P_{\text{interp}} = \frac{\sum^{K}_{k=1} P_k}{K}
\]

% demosaicado_bilineal_tikz.tex
% Gráfica TikZ del proceso de demosaicado bilineal - Versión corregida con superposición de esquinas

\begin{figure}[h]
\centering
\begin{tikzpicture}[
    font=\small,
    pixel/.style={rectangle, draw=black, minimum width=1.2cm, minimum height=1.2cm, inner sep=0}
]

% ============================================
% PATRÓN BAYER ORIGINAL (2x2 RGGB)
% ============================================

% Cuadrante superior izquierdo (rojo)
\fill[red!80!black] (-1.2,0) rectangle (0,1.2);
\draw[black, thick] (-1.2,0) rectangle (0,1.2);
\node[white, font=\bfseries] at (-0.6,0.6) (r) {$r$};

% Cuadrante superior derecho (verde 1)
\fill[green!60!black] (0,0) rectangle (1.2,1.2);
\draw[black, thick] (0,0) rectangle (1.2,1.2);
\node[white, font=\bfseries] at (0.6,0.6) (g1) {$g_1$};

% Cuadrante inferior izquierdo (verde 2)
\fill[green!40!black] (-1.2,-1.2) rectangle (0,0);
\draw[black, thick] (-1.2,-1.2) rectangle (0,0);
\node[white, font=\bfseries] at (-0.6,-0.6) (g2) {$g_2$};

% Cuadrante inferior derecho (azul)
\fill[blue!80!black] (0,-1.2) rectangle (1.2,0);
\draw[black, thick] (0,-1.2) rectangle (1.2,0);
\node[white, font=\bfseries] at (0.6,-0.6) (b) {$b$};

% Marco exterior del patrón
\draw[black, very thick] (-1.2,-1.2) rectangle (1.2,1.2);

% Título
\node[above, align=center, font=\bfseries] at (0,1.4) 
    {Patrón Bayer 2×2};

% ============================================
% FLECHA DE PROCESO
% ============================================
\node at (2.5,0) {$\Longrightarrow$};

% ============================================
% PÍXEL RESULTANTE (CARTAS SUPERPUESTAS CON ESQUINAS)
% ============================================

% Parámetros
\def\cardsize{1.2cm}
\def\overlap{0.3cm} % L/4 = 1.2cm/4 = 0.3cm
\def\shift{0.9cm} % cardsize - overlap = 1.2 - 0.3 = 0.9cm
\def\sqareorigin{3.5cm}

% ORDEN: Primero dibujamos la de atrás (B), luego la del medio (G), luego la del frente (R)

% Azul (B) - atrás - más desplazada hacia abajo y derecha
\fill[blue!80!black]                (\sqareorigin+2*\shift            , -2*\shift) rectangle (\sqareorigin+2*\shift+\cardsize, -2*\shift+\cardsize);
\draw[black, thick]                 (\sqareorigin+2*\shift            , -2*\shift) rectangle (\sqareorigin+2*\shift+\cardsize, -2*\shift+\cardsize);
\node[white, font=\bfseries] at     (\sqareorigin+2*\shift+\cardsize/2, -2*\shift+\cardsize/2) (B) {$b$};

% Verde (G) - medio - desplazada un shift
\fill[green!50!black]               (\sqareorigin+\shift              , -  \shift) rectangle (\sqareorigin+\shift+\cardsize, -\shift+\cardsize);
\draw[black, thick]                 (\sqareorigin+\shift              , -  \shift) rectangle (\sqareorigin+\shift+\cardsize, -\shift+\cardsize);
\node[white, font=\bfseries] at     (\sqareorigin+\shift+\cardsize/2  , -  \shift+\cardsize/2) (G) {$\frac{g_1 + g_2}{2}$};

% Rojo (R) - frente - posición base (sin desplazar)
\fill[red!80!black]                 (\sqareorigin,0) rectangle (\sqareorigin+\cardsize           , \cardsize);
\draw[black, thick]                 (\sqareorigin,0) rectangle (\sqareorigin+\cardsize           , \cardsize);
\node[white, font=\bfseries] at     (\sqareorigin+\cardsize/2         , \cardsize/2) (R) {$r$};

% Título del píxel resultante
\node[above, font=\bfseries] at (\sqareorigin+\cardsize/2+\shift, 1.5) {Píxel RGB};

\end{tikzpicture}
\caption{Proceso de demosaicado bilineal para patrón Bayer RGGB. Los canales de color se superponen para formar el píxel RGB final.}
\label{fig:demosaicado_bilineal}
\end{figure}


\subsection{Interpolación Bicuadrada}
Utiliza tres puntos para ajustar una parábola:

\[
P(x) = ax^2 + bx + c
\]

Para puntos equidistantes en $x \in \{-2, 0, 2\}$ (separados por 2 unidades, como en el patrón Bayer), la fórmula de interpolación para el punto intermedio en $x=1$ es:

\[
P(1) = \frac{3}{8}P(-2) + \frac{3}{4}P(0) - \frac{1}{8}P(2)
\]

En coordenadas de imagen, para el canal rojo en posición $(2k,j)$:

\[
R_{2k,j} = \frac{3}{8}R_{2k-2,j} + \frac{3}{4}R_{2k,j} - \frac{1}{8}R_{2k+2,j}
\]

% separacion_canales_bayer_tikz.tex
% Gráfica TikZ de la separación de canales en patrón Bayer - Versión exacta

\begin{figure}[h]
\centering
\begin{tikzpicture}[
    font=\normalsize,
    cell/.style={rectangle, draw=black, minimum size=0.9cm, inner sep=0}
]

% ============================================
% PATRÓN BAYER ORIGINAL 3×3
% ============================================

% Título
\node[above, font=\bfseries] at (0, 2.5) {Patrón Bayer Original};

% Fila 1: R G R
\node[cell, fill=red!80!black] at (0, 2) {\textcolor{white}{R}};
\node[cell, fill=green!60!black] at (0.9, 2) {\textcolor{white}{G}};
\node[cell, fill=red!80!black] at (1.8, 2) {\textcolor{white}{R}};

% Fila 2: G B G
\node[cell, fill=green!40!black] at (0, 1.1) {\textcolor{white}{G}};
\node[cell, fill=blue!80!black] at (0.9, 1.1) {\textcolor{white}{B}};
\node[cell, fill=green!40!black] at (1.8, 1.1) {\textcolor{white}{G}};

% Fila 3: R G R
\node[cell, fill=red!80!black] at (0, 0.2) {\textcolor{white}{R}};
\node[cell, fill=green!60!black] at (0.9, 0.2) {\textcolor{white}{G}};
\node[cell, fill=red!80!black] at (1.8, 0.2) {\textcolor{white}{R}};

% ============================================
% FLECHA DE SEPARACIÓN
% ============================================
\node at (2.7, 1.1) {$\Rightarrow$};

% ============================================
% CANAL ROJO (R)
% ============================================

% Título
\node[above, font=\bfseries, red!80!black] at (4, 2.5) {Canal R};

% Fila 1: R 0 R
\node[cell, fill=red!80!black] at (3.5, 2) {\textcolor{white}{R}};
\node[cell, fill=gray!20] at (4.4, 2) {0};
\node[cell, fill=red!80!black] at (5.3, 2) {\textcolor{white}{R}};

% Fila 2: 0 0 0
\node[cell, fill=gray!20] at (3.5, 1.1) {0};
\node[cell, fill=gray!20] at (4.4, 1.1) {0};
\node[cell, fill=gray!20] at (5.3, 1.1) {0};

% Fila 3: R 0 R
\node[cell, fill=red!80!black] at (3.5, 0.2) {\textcolor{white}{R}};
\node[cell, fill=gray!20] at (4.4, 0.2) {0};
\node[cell, fill=red!80!black] at (5.3, 0.2) {\textcolor{white}{R}};

% ============================================
% CANAL VERDE (G)
% ============================================

% Título
\node[above, font=\bfseries, green!60!black] at (7, 2.5) {Canal G};

% Fila 1: 0 G 0
\node[cell, fill=gray!20] at (6.5, 2) {0};
\node[cell, fill=green!60!black] at (7.4, 2) {\textcolor{white}{G}};
\node[cell, fill=gray!20] at (8.3, 2) {0};

% Fila 2: G 0 G
\node[cell, fill=green!40!black] at (6.5, 1.1) {\textcolor{white}{G}};
\node[cell, fill=gray!20] at (7.4, 1.1) {0};
\node[cell, fill=green!40!black] at (8.3, 1.1) {\textcolor{white}{G}};

% Fila 3: 0 G 0
\node[cell, fill=gray!20] at (6.5, 0.2) {0};
\node[cell, fill=green!60!black] at (7.4, 0.2) {\textcolor{white}{G}};
\node[cell, fill=gray!20] at (8.3, 0.2) {0};

% ============================================
% CANAL AZUL (B)
% ============================================

% Título
\node[above, font=\bfseries, blue!80!black] at (10, 2.5) {Canal B};

% Fila 1: 0 0 0
\node[cell, fill=gray!20] at (9.5, 2) {0};
\node[cell, fill=gray!20] at (10.4, 2) {0};
\node[cell, fill=gray!20] at (11.3, 2) {0};

% Fila 2: 0 B 0
\node[cell, fill=gray!20] at (9.5, 1.1) {0};
\node[cell, fill=blue!80!black] at (10.4, 1.1) {\textcolor{white}{B}};
\node[cell, fill=gray!20] at (11.3, 1.1) {0};

% Fila 3: 0 0 0
\node[cell, fill=gray!20] at (9.5, 0.2) {0};
\node[cell, fill=gray!20] at (10.4, 0.2) {0};
\node[cell, fill=gray!20] at (11.3, 0.2) {0};

\end{tikzpicture}
\caption{Separación de canales en patrón Bayer 3×3. El patrón original se descompone en tres matrices que muestran la información de cada canal por separado. Los ceros (gris) indican ausencia de información en esa posición para el canal correspondiente.}
\label{fig:separacion_canales_bayer}
\end{figure}

% interpolacion_bicuadrada_R_completo_tikz.tex
% Figura completa de interpolación bicuadrada para canal R (fórmulas con espaciado uniforme)

\begin{figure}[H]
\centering
\begin{tikzpicture}[
    font=\normalsize,
    cell/.style={rectangle, draw=black, minimum size=1.2cm, inner sep=0},
    legendcell/.style={rectangle, draw=black, minimum size=0.8cm, inner sep=0},
    formula/.style={font=\footnotesize, align=left}
]

% ============================================
% PARÁMETROS DE LA CUADRÍCULA
% ============================================
\def\gridWidth{14}   % Ancho total aumentado (en cm)
\def\rowOneHeight{3.5}  % Alto de la fila 1 (en cm)
\def\rowTwoHeight{3}    % Alto de la fila 2 (en cm)
\def\colWidth{2.8}    % Ancho de cada columna (en cm)

% ============================================
% DEFINICIÓN DE COORDENADAS DE CENTRO
% ============================================
% Fila 1 (ahora con 5 columnas para más espacio)
\coordinate (F1C1) at (0.5*\colWidth, 0);
\coordinate (F1C2) at (1.5*\colWidth, 0);
\coordinate (F1C3) at (2.5*\colWidth, 0);
\coordinate (F1C4) at (3.5*\colWidth, 0); % Espacio extra
\coordinate (F1C5) at (4.5*\colWidth, 0); % Leyenda más a la derecha

% Fila 2 (ocupa ancho completo de 5 columnas)
\coordinate (F2) at (2.5*\colWidth, -\rowOneHeight-0.5*\rowTwoHeight);

% ============================================
% F1C1: MATRIZ ORIGINAL - CANAL R
% ============================================
\begin{scope}[shift={(F1C1)}]
    % Título
    \node[above, font=\bfseries, red!80!black] at (0, 1.8) {Canal R Original};
    
    % Matriz 3x3
    % Fila 1
    \node[cell, fill=red!80!black] at (-1.2, 1.2) {\textcolor{white}{$r_{11}$}};
    \node[cell, fill=gray!20] at (0, 1.2) {$0$};
    \node[cell, fill=red!80!black] at (1.2, 1.2) {\textcolor{white}{$r_{13}$}};
    
    % Fila 2
    \node[cell, fill=gray!20] at (-1.2, 0) {$0$};
    \node[cell, fill=gray!20] at (0, 0) {$0$};
    \node[cell, fill=gray!20] at (1.2, 0) {$0$};
    
    % Fila 3
    \node[cell, fill=red!80!black] at (-1.2, -1.2) {\textcolor{white}{$r_{31}$}};
    \node[cell, fill=gray!20] at (0, -1.2) {$0$};
    \node[cell, fill=red!80!black] at (1.2, -1.2) {\textcolor{white}{$r_{33}$}};
\end{scope}

% ============================================
% F1C2: FLECHA DE PROCESO
% ============================================
\begin{scope}[shift={(F1C2)}]
    \node[font=\Huge] at (0, 0) {$\Rightarrow$};
    \node[above, font=\scriptsize] at (0, 1) {Interpolación};
    \node[below, font=\scriptsize] at (0, -1) {Bicuadrada};
\end{scope}

% ============================================
% F1C3: MATRIZ INTERPOLADA - CANAL R
% ============================================
\begin{scope}[shift={(F1C3)}]
    % Título
    \node[above, font=\bfseries, red!80!black] at (0, 1.8) {Canal R Interpolado};
    
    % Matriz 3x3 completa
    % Fila 1
    \node[cell, fill=red!80!black] at (-1.2, 1.2) {\textcolor{white}{$r_{11}$}};
    \node[cell, fill=red!60!black] at (0, 1.2) {\textcolor{white}{$r_{12}$}};
    \node[cell, fill=red!80!black] at (1.2, 1.2) {\textcolor{white}{$r_{13}$}};
    
    % Fila 2
    \node[cell, fill=red!60!black] at (-1.2, 0) {\textcolor{white}{$r_{21}$}};
    \node[cell, fill=red!40!black] at (0, 0) {\textcolor{white}{$r_{22}$}};
    \node[cell, fill=red!60!black] at (1.2, 0) {\textcolor{white}{$r_{23}$}};
    
    % Fila 3
    \node[cell, fill=red!80!black] at (-1.2, -1.2) {\textcolor{white}{$r_{31}$}};
    \node[cell, fill=red!60!black] at (0, -1.2) {\textcolor{white}{$r_{32}$}};
    \node[cell, fill=red!80!black] at (1.2, -1.2) {\textcolor{white}{$r_{33}$}};
\end{scope}

% ============================================
% F1C5: LEYENDA DE COLORES (AJUSTADA VERTICALMENTE)
% ============================================
\begin{scope}[shift={(F1C5)}]
    % Título
    \node[above, font=\bfseries] at (0, 1.8) {Leyenda de Colores};
    
    % Ajuste: movemos toda la leyenda 0.8cm hacia arriba
    \def\verticalShift{0.8}
    
    % Original (rojo oscuro)
    \node[legendcell, fill=red!80!black] at (-1, 0.5+\verticalShift) {};
    \node[anchor=west, font=\small] at (-0.7, 0.5+\verticalShift) {Original};
    \node[anchor=west, font=\tiny] at (-0.7, 0.2+\verticalShift) {($r_{11}$, $r_{13}$, $r_{31}$, $r_{33}$)};
    
    % Interpolado 1D (rojo medio)
    \node[legendcell, fill=red!60!black] at (-1, -0.3+\verticalShift) {};
    \node[anchor=west, font=\small] at (-0.7, -0.3+\verticalShift) {Interpolado 1D};
    \node[anchor=west, font=\tiny] at (-0.7, -0.6+\verticalShift) {($r_{12}$, $r_{21}$, $r_{23}$, $r_{32}$)};
    
    % Interpolado 2D (rojo claro)
    \node[legendcell, fill=red!40!black] at (-1, -1.1+\verticalShift) {};
    \node[anchor=west, font=\small] at (-0.7, -1.1+\verticalShift) {Interpolado 2D};
    \node[anchor=west, font=\tiny] at (-0.7, -1.4+\verticalShift) {($r_{22}$)};
    
    % Por interpolar (gris)
    \node[legendcell, fill=gray!20] at (-1, -1.9+\verticalShift) {};
    \node[anchor=west, font=\small] at (-0.7, -1.9+\verticalShift) {Por interpolar};
    \node[anchor=west, font=\tiny] at (-0.7, -2.2+\verticalShift) {(valor 0)};
\end{scope}

% ============================================
% F2: FÓRMULAS DE INTERPOLACIÓN CON ESPACIADO UNIFORME
% ============================================
\begin{scope}[shift={(F2)}]
    % Título
    \node[above, font=\bfseries] at (0, 1.5) {Fórmulas de Interpolación Bicuadrada};
    
    % Parámetros para el espaciado uniforme
    \def\formulaStartY{0.8}      % Coordenada Y inicial
    \def\formulaSpacing{0.7}     % Espaciado uniforme entre fórmulas
    \def\leftColumnX{-4}         % Columna izquierda
    \def\rightColumnX{2}         % Columna derecha (separada 6cm de la izquierda)
    
    % Columna izquierda: 3 fórmulas
    \node[anchor=north west, formula] at (\leftColumnX, \formulaStartY) 
        {$r_{12} = \tfrac{3}{8}r_{11} + \tfrac{3}{4}r_{13} - \tfrac{1}{8}r_{15}$};
    
    \node[anchor=north west, formula] at (\leftColumnX, \formulaStartY - \formulaSpacing) 
        {$r_{21} = \tfrac{3}{8}r_{11} + \tfrac{3}{4}r_{31} - \tfrac{1}{8}r_{51}$};
    
    \node[anchor=north west, formula] at (\leftColumnX, \formulaStartY - 2*\formulaSpacing) 
        {$r_{22} = \tfrac{9}{64}r_{11} + \tfrac{9}{32}r_{13} + \tfrac{9}{32}r_{31} + \tfrac{9}{16}r_{33} - \tfrac{3}{64}r_{15} - \tfrac{3}{64}r_{51}$};
    
    % Columna derecha: 2 fórmulas (alineadas con las dos primeras de la izquierda)
    \node[anchor=north west, formula] at (\rightColumnX, \formulaStartY) 
        {$r_{23} = \tfrac{3}{8}r_{13} + \tfrac{3}{4}r_{33} - \tfrac{1}{8}r_{53}$};
    
    \node[anchor=north west, formula] at (\rightColumnX, \formulaStartY - \formulaSpacing) 
        {$r_{32} = \tfrac{3}{8}r_{31} + \tfrac{3}{4}r_{33} - \tfrac{1}{8}r_{35}$};
    
    % Línea vertical divisoria (opcional, para referencia)
    \draw[gray!30, dashed] (\leftColumnX + 3, \formulaStartY + 0.2) -- (\leftColumnX + 3, \formulaStartY - 2*\formulaSpacing - 0.2);
\end{scope}

% ============================================
% LÍNEA DIVISORIA HORIZONTAL
% ============================================
\draw[gray!40, thick] (0.2, -\rowOneHeight+0.2) -- (\gridWidth-0.2, -\rowOneHeight+0.2);

\end{tikzpicture}
\caption{Interpolación bicuadrada para el canal rojo (R) en patrón Bayer. Arriba se muestra la transformación de la matriz original (con valores conocidos en rojo oscuro y ceros en gris) a la matriz interpolada (todos los valores completos, diferenciando entre originales, interpolados 1D e interpolados 2D). La leyenda, alineada verticalmente con las matrices, explica los colores utilizados. Debajo se presentan las fórmulas de interpolación bicuadrada con espaciado uniforme, distribuidas en dos columnas para una presentación clara y ordenada.}
\label{fig:interpolacion_bicuadrada_R_completa}
\end{figure}


\subsection{Interpolación Bicúbica}
Utiliza cuatro puntos para ajustar un polinomio cúbico:

\[
P(x) = ax^3 + bx^2 + cx + d
\]

Con puntos en $x \in \{-2, 0, 2, 4\}$, usando el kernel cúbico de Catmull-Rom:

\[
P(1) = -\frac{1}{16}P(-2) + \frac{9}{16}P(0) + \frac{9}{16}P(2) - \frac{1}{16}P(4)
\]

Para el canal verde en posición $(2k+1,2\ell)$:

\[
G_{2k+1,2\ell} = -\frac{1}{16}G_{2k-3,2\ell} + \frac{9}{16}G_{2k-1,2\ell} + \frac{9}{16}G_{2k+3,2\ell} - \frac{1}{16}G_{2k+5,2\ell}
\]

% interpolacion_bicubica_R_completo_tikz_final5.tex
% Versión final con fórmulas más a la izquierda y títulos separados

\begin{figure}[h]
\centering
\begin{tikzpicture}[
    font=\normalsize,
    cell/.style={rectangle, draw=black, minimum size=0.9cm, inner sep=0},
    legendcell/.style={rectangle, draw=black, minimum size=0.8cm, inner sep=0},
    formula/.style={font=\footnotesize, align=left}
]

% ============================================
% PARÁMETROS DE LA CUADRÍCULA
% ============================================
\def\gridWidth{16}   % Ancho total aumentado (en cm)
\def\rowOneHeight{3.5}  % Alto de la fila 1 (en cm)
\def\rowTwoHeight{5.0}  % Alto de la fila 2 reducido (en cm)
\def\colWidth{2.8}    % Ancho de cada columna (en cm)

% ============================================
% DEFINICIÓN DE COORDENADAS DE CENTRO
% ============================================
% Fila 1 (5 columnas)
\coordinate (F1C1) at (0.5*\colWidth, 0);
\coordinate (F1C2) at (1.5*\colWidth, 0);
\coordinate (F1C3) at (2.5*\colWidth, 0);
\coordinate (F1C4) at (3.5*\colWidth, 0);
\coordinate (F1C5) at (4.5*\colWidth, 0);

% Fila 2 (ocupa ancho completo)
\coordinate (F2) at (2.5*\colWidth, -\rowOneHeight-0.5*\rowTwoHeight);

% ============================================
% F1C1: MATRIZ ORIGINAL - CANAL R (4×4) CORREGIDA
% ============================================
\begin{scope}[shift={(F1C1)}]
    % Título
    \node[above, font=\bfseries, red!80!black] at (0, 2.2) {Canal R Original (4×4)};
    
    % Matriz 4x4 con valores en posiciones correctas según patrón Bayer
    % Fila 1 (índice 1)
    \node[cell, fill=red!80!black] at (-1.35, 1.35) {\textcolor{white}{$r_{11}$}};
    \node[cell, fill=gray!20] at (-0.45, 1.35) {$0$};
    \node[cell, fill=red!80!black] at (0.45, 1.35) {\textcolor{white}{$r_{13}$}};
    \node[cell, fill=gray!20] at (1.35, 1.35) {$0$};
    
    % Fila 2 (índice 2)
    \node[cell, fill=gray!20] at (-1.35, 0.45) {$0$};
    \node[cell, fill=gray!20] at (-0.45, 0.45) {$0$};
    \node[cell, fill=gray!20] at (0.45, 0.45) {$0$};
    \node[cell, fill=gray!20] at (1.35, 0.45) {$0$};
    
    % Fila 3 (índice 3)
    \node[cell, fill=red!80!black] at (-1.35, -0.45) {\textcolor{white}{$r_{31}$}};
    \node[cell, fill=gray!20] at (-0.45, -0.45) {$0$};
    \node[cell, fill=red!80!black] at (0.45, -0.45) {\textcolor{white}{$r_{33}$}};
    \node[cell, fill=gray!20] at (1.35, -0.45) {$0$};
    
    % Fila 4 (índice 4)
    \node[cell, fill=gray!20] at (-1.35, -1.35) {$0$};
    \node[cell, fill=gray!20] at (-0.45, -1.35) {$0$};
    \node[cell, fill=gray!20] at (0.45, -1.35) {$0$};
    \node[cell, fill=gray!20] at (1.35, -1.35) {$0$};
\end{scope}

% ============================================
% F1C2: FLECHA DE PROCESO
% ============================================
\begin{scope}[shift={(F1C2)}]
    \node[font=\Huge] at (0, 0) {$\Rightarrow$};
    \node[above, font=\scriptsize] at (0, 1) {Interpolación};
    \node[below, font=\scriptsize] at (0, -1) {Bicúbica};
\end{scope}

% ============================================
% F1C3: MATRIZ INTERPOLADA - CANAL R (4×4)
% ============================================
\begin{scope}[shift={(F1C3)}]
    % Título
    \node[above, font=\bfseries, red!80!black] at (0, 2.2) {Canal R Interpolado (4×4)};
    
    % Matriz 4x4 completa
    % Fila 1
    \node[cell, fill=red!80!black] at (-1.35, 1.35) {\textcolor{white}{$r_{11}$}};
    \node[cell, fill=red!70!black] at (-0.45, 1.35) {\textcolor{white}{$r_{12}$}};
    \node[cell, fill=red!80!black] at (0.45, 1.35) {\textcolor{white}{$r_{13}$}};
    \node[cell, fill=red!70!black] at (1.35, 1.35) {\textcolor{white}{$r_{14}$}};
    
    % Fila 2
    \node[cell, fill=red!70!black] at (-1.35, 0.45) {\textcolor{white}{$r_{21}$}};
    \node[cell, fill=red!50!black] at (-0.45, 0.45) {\textcolor{white}{$r_{22}$}};
    \node[cell, fill=red!60!black] at (0.45, 0.45) {\textcolor{white}{$r_{23}$}};
    \node[cell, fill=red!70!black] at (1.35, 0.45) {\textcolor{white}{$r_{24}$}};
    
    % Fila 3
    \node[cell, fill=red!80!black] at (-1.35, -0.45) {\textcolor{white}{$r_{31}$}};
    \node[cell, fill=red!60!black] at (-0.45, -0.45) {\textcolor{white}{$r_{32}$}};
    \node[cell, fill=red!80!black] at (0.45, -0.45) {\textcolor{white}{$r_{33}$}};
    \node[cell, fill=red!60!black] at (1.35, -0.45) {\textcolor{white}{$r_{34}$}};
    
    % Fila 4
    \node[cell, fill=red!70!black] at (-1.35, -1.35) {\textcolor{white}{$r_{41}$}};
    \node[cell, fill=red!50!black] at (-0.45, -1.35) {\textcolor{white}{$r_{42}$}};
    \node[cell, fill=red!60!black] at (0.45, -1.35) {\textcolor{white}{$r_{43}$}};
    \node[cell, fill=red!70!black] at (1.35, -1.35) {\textcolor{white}{$r_{44}$}};
\end{scope}

% ============================================
% F1C5: LEYENDA DE COLORES CORREGIDA
% ============================================
\begin{scope}[shift={(F1C5)}]
    % Título
    \node[above, font=\bfseries] at (0, 2.2) {Leyenda de Colores};
    
    % Ajuste vertical
    \def\verticalShift{0.6}
    
    % Original (rojo oscuro)
    \node[legendcell, fill=red!80!black] at (-1, 0.8+\verticalShift) {};
    \node[anchor=west, font=\small] at (-0.7, 0.8+\verticalShift) {Original};
    \node[anchor=west, font=\tiny] at (-0.7, 0.5+\verticalShift) {($r_{11}$, $r_{13}$, $r_{31}$, $r_{33}$)};
    
    % Interpolado 1D (rojo medio-alto)
    \node[legendcell, fill=red!70!black] at (-1, -0.1+\verticalShift) {};
    \node[anchor=west, font=\small] at (-0.7, -0.1+\verticalShift) {Interp. Borde};
    \node[anchor=west, font=\tiny] at (-0.7, -0.4+\verticalShift) {($r_{12}$, $r_{14}$, $r_{21}$, $r_{24}$)};
    
    % Interpolado 2D (rojo medio)
    \node[legendcell, fill=red!50!black] at (-1, -1.0+\verticalShift) {};
    \node[anchor=west, font=\small] at (-0.7, -1.0+\verticalShift) {Interp. 2D};
    \node[anchor=west, font=\tiny] at (-0.7, -1.3+\verticalShift) {($r_{22}$, $r_{23}$, $r_{32}$, $r_{34}$)};
    
    % Interpolado centro (rojo claro)
    \node[legendcell, fill=red!30!black] at (-1, -1.9+\verticalShift) {};
    \node[anchor=west, font=\small] at (-0.7, -1.9+\verticalShift) {Interp. Centro};
    \node[anchor=west, font=\tiny] at (-0.7, -2.2+\verticalShift) {($r_{42}$, $r_{43}$, $r_{44}$)};
    
    % Por interpolar (gris)
    \node[legendcell, fill=gray!20] at (-1, -2.8+\verticalShift) {};
    \node[anchor=west, font=\small] at (-0.7, -2.8+\verticalShift) {Por interpolar};
    \node[anchor=west, font=\tiny] at (-0.7, -3.1+\verticalShift) {(valor 0)};
\end{scope}

% ============================================
% F2: FÓRMULAS AÚN MÁS A LA IZQUIERDA CON TÍTULOS SEPARADOS
% ============================================
\begin{scope}[shift={(F2)}]
    % Título principal
    \node[above, font=\bfseries] at (0, 1.7) {Fórmulas de Interpolación Bicúbica};
    
    % Parámetros para espaciado uniforme - ajustados
    \def\titleSpacing{0.6}       % ESPACIO AUMENTADO después del título de sección
    \def\formulaSpacing{0.65}    % Espacio uniforme entre fórmulas
    \def\sectionSpacing{0.6}     % Espacio entre secciones
    \def\columnX{-6.5}           % COLUMNA MUCHO MÁS A LA IZQUIERDA
    
    % Coordenadas Y fijas - ajustadas
    \def\startY{1.5}
    \def\yOneTitle{\startY}
    \def\yOneA{\startY - \titleSpacing}        % Más espacio después del título
    \def\yOneB{\yOneA - \formulaSpacing}
    \def\yOneC{\yOneB - \formulaSpacing}
    \def\yOneD{\yOneC - \formulaSpacing}
    
    \def\yTwoTitle{\yOneD - \sectionSpacing}
    \def\yTwoA{\yTwoTitle - \titleSpacing}     % Más espacio después del título
    \def\yTwoB{\yTwoA - \formulaSpacing}
    \def\yTwoC{\yTwoB - \formulaSpacing}
    \def\yTwoD{\yTwoC - \formulaSpacing}
    
    % ========== INTERPOLACIONES 1D ==========
    % Subtítulo: Interpolación 1D - con más espacio debajo
    \node[anchor=west, font=\small\bfseries, align=left] at (\columnX, \yOneTitle) 
        {Interpolación 1D:};
    
    % Fórmulas 1D (alineadas a la izquierda)
    \node[anchor=west, formula, align=left] at (\columnX, \yOneA) 
        {$r_{12} = -\tfrac{1}{16}r_{10} + \tfrac{9}{16}r_{11} + \tfrac{9}{16}r_{13} - \tfrac{1}{16}r_{14}$};
    
    \node[anchor=west, formula, align=left] at (\columnX, \yOneB) 
        {$r_{14} = -\tfrac{1}{16}r_{12} + \tfrac{9}{16}r_{13} + \tfrac{9}{16}r_{15} - \tfrac{1}{16}r_{16}$};
    
    \node[anchor=west, formula, align=left] at (\columnX, \yOneC) 
        {$r_{21} = -\tfrac{1}{16}r_{01} + \tfrac{9}{16}r_{11} + \tfrac{9}{16}r_{31} - \tfrac{1}{16}r_{41}$};
    
    \node[anchor=west, formula, align=left] at (\columnX, \yOneD) 
        {$r_{41} = -\tfrac{1}{16}r_{21} + \tfrac{9}{16}r_{31} + \tfrac{9}{16}r_{51} - \tfrac{1}{16}r_{61}$};
    
    % ========== INTERPOLACIONES 2D ==========
    % Subtítulo: Interpolación 2D - con más espacio debajo
    \node[anchor=west, font=\small\bfseries, align=left] at (\columnX, \yTwoTitle) 
        {Interpolación 2D:};
    
    % Fórmulas 2D (alineadas a la izquierda)
    \node[anchor=west, formula, align=left, text width=14.5cm] at (\columnX, \yTwoA) 
        {$r_{22} = \tfrac{1}{256}(r_{00} - 9r_{01} - 9r_{10} + 81r_{11} + 81r_{13} - 9r_{14} - 9r_{30} + 81r_{31} + 81r_{33} - 9r_{34})$};
    
    \node[anchor=west, formula, align=left, text width=14.5cm] at (\columnX, \yTwoB) 
        {$r_{23} = \tfrac{1}{256}(r_{01} - 9r_{02} - 9r_{11} + 81r_{12} + 81r_{14} - 9r_{15} - 9r_{31} + 81r_{32} + 81r_{34} - 9r_{35})$};
    
    \node[anchor=west, formula, align=left, text width=14.5cm] at (\columnX, \yTwoC) 
        {$r_{32} = \tfrac{1}{256}(r_{10} - 9r_{11} - 9r_{20} + 81r_{21} + 81r_{23} - 9r_{24} - 9r_{40} + 81r_{41} + 81r_{43} - 9r_{44})$};
    
    \node[anchor=west, formula, align=left, text width=14.5cm] at (\columnX, \yTwoD) 
        {$r_{33} = \tfrac{1}{256}(r_{11} - 9r_{12} - 9r_{21} + 81r_{22} + 81r_{24} - 9r_{25} - 9r_{41} + 81r_{42} + 81r_{44} - 9r_{45})$};
\end{scope}

% ============================================
% LÍNEA DIVISORIA HORIZONTAL
% ============================================
\draw[gray!40, thick] (0.2, -\rowOneHeight+0.2) -- (\gridWidth-0.2, -\rowOneHeight+0.2);

\end{tikzpicture}
\caption{Interpolación bicúbica para el canal rojo (R) en patrón Bayer 4×4. Arriba se muestran las matrices original e interpolada, con diferentes tonos de rojo indicando el nivel de interpolación. La leyenda explica la codificación de colores. Debajo se presentan las fórmulas organizadas verticalmente: interpolaciones 1D (para bordes) y 2D (para regiones internas) con el kernel bicúbico completo que utiliza un vecindario de 4×4 píxeles.}
\label{fig:interpolacion_bicubica_R_completa_final5}
\end{figure}


\subsection{Proceso Bidimensional}
La interpolación 2D se realiza en dos pasos independientes: primero se interpola verticalmente y luego horizontalmente (o viceversa). Este enfoque separable reduce la complejidad computacional respecto a una interpolación 2D directa.
% resultados.tex
\section{Resultados}\label{sec:resultados}

\subsection{Implementación Práctica}

Los algoritmos fueron implementados en MATLAB y probados con imágenes sintéticas y reales. 
Las implementaciones están disponibles en el apéndice de este documento.

\subsection{Complejidad Computacional}

\begin{table}[h]
\centering
\caption{Comparación de complejidad de métodos}
\begin{tabular}{lccc}
\hline
\textbf{Método} & \textbf{Operaciones por pixel} & \textbf{Memoria} & \textbf{Tiempo relativo} \\
\hline
Lineal & 2-4 & Baja & 1.0× \\
Bicuadrada & 8-12 & Media & 2.5× \\
Bicúbica & 16-24 & Alta & 4.0× \\
\hline
\end{tabular}
\end{table}

\subsection{Calidad de Interpolación}

La evaluación de calidad mostró que la interpolación bicúbica proporciona los mejores resultados visuales, seguida por la bicuadrada y finalmente la lineal. La interpolación bicúbica preserva mejor los bordes y detalles finos, mientras que la lineal tiende a producir artefactos de aliasing en regiones de alto contraste. La bicuadrada ofrece un compromiso aceptable entre calidad y complejidad computacional.

\subsection{Problemas en Bordes}

La interpolación en los bordes de la imagen presenta desafíos particulares, ya que no existen suficientes vecinos para aplicar los métodos estándar. Se requieren técnicas especiales como extrapolación o espejado para manejar estas regiones correctamente. Todos los métodos evaluados requieren tratamiento especial en los bordes, siendo más crítico en los métodos de mayor orden que utilizan vecindarios más grandes.

\subsection{Vecindarios Utilizados}

\begin{table}[h]
\centering
\caption{Configuraciones de vecindario para interpolación}
\begin{tabular}{lp{3cm}p{3cm}p{3cm}}
\hline
\textbf{Método} & \textbf{Patrón 1D} & \textbf{Tamaño} & \textbf{Puntos utilizados} \\
\hline
Lineal & [P₁ ? P₂] & 3×1 & 2 \\
Bicuadrada & [P₁ P₂ ? P₃ P₄] & 5×1 & 3 \\
Bicúbica & [P₁ P₂ P₃ ? P₄ P₅ P₆] & 7×1 & 4 \\
\hline
\end{tabular}
\end{table}
% discusion.tex
\section{Discusión}\label{sec:discusion}

\subsection{Comparación de Métodos}

Los tres métodos de interpolación presentan diferentes compromisos entre calidad y complejidad:

\subsubsection{Interpolación Lineal}
\begin{itemize}
    \item \textbf{Ventajas:} Simplicidad computacional, baja latencia, implementación trivial
    \item \textbf{Desventajas:} Artefactos de aliasing en bordes, pérdida de detalle fino
    \item \textbf{Aplicaciones:} Sistemas en tiempo real con recursos limitados
\end{itemize}

\subsubsection{Interpolación Bicuadrada}
\begin{itemize}
    \item \textbf{Ventajas:} Mejor calidad que lineal, suavizado moderado, complejidad aceptable
    \item \textbf{Desventajas:} Puede suavizar en exceso, requiere vecindario de 5 puntos
    \item \textbf{Aplicaciones:} Fotografía de consumo, video de calidad media
\end{itemize}

\subsubsection{Interpolación Bicúbica}
\begin{itemize}
    \item \textbf{Ventajas:} Alta calidad, preservación de bordes, resultados visualmente superiores
    \item \textbf{Desventajas:} Alta complejidad computacional, mayor latencia
    \item \textbf{Aplicaciones:} Fotografía profesional, procesamiento offline, aplicaciones médicas
\end{itemize}

\subsection{Consideraciones Prácticas}

La selección del método de interpolación debe considerar tanto los requisitos de calidad como las limitaciones computacionales del sistema. Para aplicaciones en tiempo real con recursos limitados, la interpolación lineal puede ser suficiente. En sistemas donde la calidad visual es prioritaria, los métodos bicuadrados o bicúbicos son preferibles, aunque requieren mayor capacidad de procesamiento.

\subsection{Limitaciones y Mejoras}

\subsubsection{Limitaciones Identificadas}
\begin{enumerate}
    \item \textbf{Artefactos de color:} Aparecen en regiones de alto contraste
    \item \textbf{Pérdida de resolución:} La interpolación nunca recupera información perdida
    \item \textbf{Dependencia del patrón:} Los resultados varían con diferentes patrones Bayer
    \item \textbf{Complejidad en bordes:} Requieren tratamiento especial
\end{enumerate}

\subsubsection{Mejoras Posibles}
\begin{itemize}
    \item \textbf{Adaptabilidad:} Usar diferentes métodos según el contenido local
    \item \textbf{Corrección de color:} Post-procesamiento para reducir artefactos
    \item \textbf{Interpolación direccional:} Considerar la orientación de bordes
    \item \textbf{Fusión de múltiples métodos:} Combinar ventajas de diferentes enfoques
\end{itemize}
% conclusion.tex
\section{Conclusión}\label{sec:conclusion}

\subsection{Conclusiones Principales}

Este estudio implementó y comparó tres métodos de interpolación para el patrón Bayer RGGB:

1. **Interpolación lineal:** El método más simple y rápido, adecuado para sistemas con recursos limitados.
2. **Interpolación bicuadrada:** Ofrece un buen equilibrio entre calidad y complejidad computacional.
3. **Interpolación bicúbica:** Proporciona la mejor calidad visual a costa de mayor complejidad computacional.

Los resultados muestran que cada método tiene su ámbito de aplicación ideal, dependiendo de los requisitos de calidad y las restricciones computacionales del sistema.

\subsection{Contribuciones del Trabajo}

\begin{enumerate}
    \item \textbf{Implementación educativa:} Código MATLAB documentado y modular
    \item \textbf{Análisis comparativo:} Evaluación sistemática de métodos polinomiales
    \item \textbf{Visualización didáctica:} Diagramas explicativos del proceso matemático
    \item \textbf{Documentación completa:} Explicación paso a paso de algoritmos
\end{enumerate}

\subsection{Recomendaciones de Uso}

\begin{table}[h]
\centering
\caption{Guía de selección de método de interpolación}
\begin{tabular}{p{4cm}p{6cm}}
\hline
\textbf{Escenario de aplicación} & \textbf{Método recomendado} \\
\hline
Sistemas embebidos en tiempo real & Lineal \\
Fotografía de consumo & Bicuadrada \\
Video de alta calidad & Bicuadrada adaptativa \\
Fotografía profesional & Bicúbica \\
Procesamiento médico/científico & Bicúbica con post-procesamiento \\
Prototipado rápido & Lineal o Bicuadrada \\
\hline
\end{tabular}
\end{table}

\subsection{Trabajo Futuro}

\begin{itemize}
    \item \textbf{Interpolación adaptativa:} Métodos que ajustan parámetros según contenido local
    \item \textbf{Aprendizaje automático:} Redes neuronales para demosaicing
    \item \textbf{Optimización hardware:} Implementación en FPGA/GPU
    \item \textbf{Patrones alternativos:} Extensiones a X-Trans y otros patrones
    \item \textbf{Corrección de artefactos:} Post-procesamiento para reducir aliasing
    \item \textbf{Benchmarking extensivo:} Comparación con métodos comerciales
\end{itemize}

\subsection{Consideraciones Finales}

La selección del método de interpolación debe basarse en un análisis cuidadoso de los requisitos 
de la aplicación específica, considerando:

\begin{itemize}
    \item Restricciones computacionales y de tiempo
    \item Calidad visual requerida
    \item Características del contenido de la imagen
    \item Recursos hardware disponibles
    \item Costo de implementación y mantenimiento
\end{itemize}

Los métodos presentados en este trabajo proporcionan una base sólida para la implementación 
de sistemas de procesamiento de imágenes que requieren interpolación de patrón Bayer, con 
diferentes niveles de compromiso entre calidad y complejidad.

% ==================== APÉNDICE: CÓDIGO ====================
\appendix
\section{Códigos Fuente}

Esta sección contiene los archivos de código implementados para este proyecto. 
Cada archivo puede descargarse haciendo clic en el enlace correspondiente.

\begin{itemize}
    \item \href{run:code/bicuadrada.m}{\texttt{bicuadrada.m}} - Implementación de interpolación bicuadrada
    \item \href{run:code/bicubica.m}{\texttt{bicubica.m}} - Implementación de interpolación bicúbica
    \item \href{run:code/comparacion.m}{\texttt{comparacion.m}} - Script de comparación de métodos
    \item \href{run:code/visualizacion.m}{\texttt{visualizacion.m}} - Visualización de resultados
\end{itemize}

\subsection{Implementación en MATLAB}

\lstinputlisting[language=Matlab, caption=Interpolación Bicuadrada (bicuadrada.m)]{./code/bicua_interp.m}

\lstinputlisting[language=Matlab, caption=Interpolación Bicúbica (bicubica.m)]{./code/bicub_interp.m}

\end{document}
